\begin{question}

  Consider a ket space spanned by the eigenkets ${\ket{a'}}$ of a
  Hermitian operator $A$.  There is no degeneracy.

  \begin{questionpart}
    Prove that
    \begin{equation*}
      \prod_{a'}(A - a')
    \end{equation*}
    is the null operator.
  \end{questionpart}

  \begin{questionpart}
    What is the significance of
    \begin{equation*}
      \prod_{a'' \ne a'}\frac{(A-a'')}{(a'-a'')}?
    \end{equation*}
  \end{questionpart}

  \begin{questionpart}
    Illustrate (a) and (b) using $A$ set equal to $S_z$ of a spin $\frac 1 2$ system.
  \end{questionpart}

\end{question}

\begin{answer}


%%%%%%%%%%%%%%%%%%%%%%%%%%%%%%%%%%%%%%%%%%%%%%%%%%%%%%%%%%%%%%%%%%%%%%%%%%%%%%%
%%                              Part a                                       %%
%%%%%%%%%%%%%%%%%%%%%%%%%%%%%%%%%%%%%%%%%%%%%%%%%%%%%%%%%%%%%%%%%%%%%%%%%%%%%%%

\begin{answerpart}{Null Operator}
  This problem will be solved in two different ways, one more elegant
  than the other.  I initially solved this problem using a direct
  combinatorial analysis (a friend of mine employed a
  proof-by-induction approach as he was most recently coming from a
  PhD in computer science).  The more elegant approach is to recognize
  the deeper significance of the individual terms and to show the
  effect.

  To begin the combinatorial approach, recite a few of the
  mathematical assumptions from the problem:

  \begin{equation}
    A^\dagger = A = \sum_i a_i\Lambda_i, \text{ where } \eval{a_i \ne a_j}_{i \ne j}
  \end{equation}

  We then need to show that

  \begin{equation}
    P \equiv \prod_i(A - a_iI) = \hat{0}
  \end{equation}

  \begin{equation}
    \begin{split}
      P
      =\ &\prod_i(A - a_iI)\\
      =\ &\prod_i(\sum_ja_j\Lambda_j - a_iI)\\
      =\ 
      &(a_1\Lambda_1 + a_2\Lambda_2 + \ldots + a_n\Lambda_n - a_1)\\
      &(a_1\Lambda_1 + a_2\Lambda_2 + \ldots + a_n\Lambda_n - a_2)\\
      &\vdots\\
      &(a_1\Lambda_1 + a_2\Lambda_2 + \ldots + a_n\Lambda_n - a_n)\\
    \end{split}
  \end{equation}

  \begin{itemize}

  \item The resulting product is all combinations of one element from each row.

  \item We know that all cross-terms including any
    $\eval{\Lambda_i\Lambda_j}_{i \ne j}$ terms will go to 0.
    That means any cross-terms with dissimilar $\Lambda$'s are 0,
    leaving only cross-terms originating from the final column.

  \item For any term of the form $\eval{a_i^k\Lambda_i}_{k>1}$,
    there will be a 1:1 correspondance with a negative term because:

    \begin{equation}
      \begin{split}
        a_i\Lambda_i \cdot a_i\Lambda_i &= a_i^2\Lambda_i\\
        -a_i \cdot a_i\Lambda_i &= -a_i^2\Lambda_i\\
      \end{split}
    \end{equation}

  \item The remaining terms are:

    \begin{equation}
      (-1)^{n-1}\sum_i\prod_ja_j\Lambda_i + (-1)^{n}\prod_ja_j
    \end{equation}

    A quick note about the $n-1$ vs $n$ for the terms there which is
    important.  In the $n-1$ case, one of the terms originates from a
    term with a $\Lambda_i$ which will always be positive, whereas all
    the other terms must originate from the right column which will
    always be negative.  Hence, there are $n-1$ negative terms applied
    and one positive, resulting in $(-1)^{n-1}$.  The final term,
    however, is the term resulting from the product of all of the bare
    terms in the final column, all $n$ of which are negative,
    resulting in $(-1)^n$.  This distinction is critical to showing
    that the two portions cancel each other out.

  \item Using $I = \sum_i\Lambda_i$ , we can show that
    
    \begin{equation}
      \begin{split}
        &(-1)^{n-1}\sum_i\prod_ja_j\Lambda_i + (-1)^{n}\prod_ja_jI\\
        =\ &(-1)^{n-1}\prod_ja_jI + (-1)^{n}\prod_ja_jI\\
        =\ &(-1)^{n-1}\prod_ja_jI - (-1)^{n-1}\prod_ja_jI\\
        =\ &(-1)^{n-1}\prod_ja_j(I - I)\\
        =\ &(-1)^{n-1}\prod_ja_j(0)\\
        =\ &0 \qed\\
      \end{split}
    \end{equation}
  \end{itemize}

  I understand that this approach is somewhat brutish, but it's always
  nice to then compare it with the elegant approach.  One of the best
  pieces of advice I received in Calc 2 was this: ``If you're ever
  stuck and don't know what to do, just start taking derivatives.''
  Similarly, this problem is a good reminder that when you're stuck,
  you should try acting a couple of kets on things to see what happens.

  \begin{equation}
    (A - a_i)\ket{a_i} = (a_i - a_i)\ket{a_i} = 0
  \end{equation}

  Each term essentially eliminates one of the eigenbases.  If the
  eigenbasis is non-degenerate, then this operator eliminates the
  entire space.$\qed$
  
\end{answerpart}


%%%%%%%%%%%%%%%%%%%%%%%%%%%%%%%%%%%%%%%%%%%%%%%%%%%%%%%%%%%%%%%%%%%%%%%%%%%%%%%
%%                              Part b                                       %%
%%%%%%%%%%%%%%%%%%%%%%%%%%%%%%%%%%%%%%%%%%%%%%%%%%%%%%%%%%%%%%%%%%%%%%%%%%%%%%%

\begin{answerpart}{Projection Operator}

  Jumping right to the easy way of doing things this time, we see that:
  
  \begin{equation}
    \begin{split}
      X\ket{a_j}
      &= \prod_{i \ne j}\frac{(A - a_i)}{(a_j - a_i)}\ket{a_j}\\
      &= \prod_{i \ne j}\frac{(a_j - a_i)}{(a_j - a_i)}\ket{a_j}\\
      &= \ket{a_j}\\
    \end{split}
  \end{equation}

  The filter on the top removes anything other than the omitted
  eigenket, and the normalization constant on the bottom ensures that
  the magnitude is 1.  Showing a negative result is also pretty easy:

  \begin{equation}
    \begin{split}
      \eval{X\ket{a_k}}_{k \ne j}
      &= \prod_{i \ne j}\frac{(A - a_i)}{(a_j - a_i)}\ket{a_k}\\
      &= \frac{A - a_k}{a_j - a_k}\prod_{i \ne j \ne k}\frac{(A - a_i)}{(a_j - a_i)}\ket{a_k}\\
      &= \frac{a_k - a_k}{a_j - a_k}\prod_{i \ne j \ne k}\frac{(A - a_i)}{(a_j - a_i)}\ket{a_k}\\
      &= 0\prod_{i \ne j \ne k}\frac{(A - a_i)}{(a_j - a_i)}\ket{a_k}\\
      &= 0\\
    \end{split}
  \end{equation}

  \begin{equation}
    \prod_{i \ne j}\frac{(A - a_i)}{(a_j - a_i)} = \Lambda_j \qed
  \end{equation}

\end{answerpart}


%%%%%%%%%%%%%%%%%%%%%%%%%%%%%%%%%%%%%%%%%%%%%%%%%%%%%%%%%%%%%%%%%%%%%%%%%%%%%%%
%%                              Part c                                       %%
%%%%%%%%%%%%%%%%%%%%%%%%%%%%%%%%%%%%%%%%%%%%%%%%%%%%%%%%%%%%%%%%%%%%%%%%%%%%%%%

\begin{answerpart}{Grind it Out for $S_z$}

  \begin{equation}
    \begin{split}
      \prod_i\left(A - a_i\right)
      &= (A - aI)(A + aI)\\
      &= AA + aA - aA - aaI\\
      &= AA + \cancel{aA} - \cancel{aA} - aaI\\
      &= AA  - aaI\\
      &= a^2(\Lambda_+ - \Lambda_-)(\Lambda_+ - \Lambda_-) - a^2I\\
      &= a^2(\Lambda_+ + \Lambda_-) - a^2I\\
      &= a^2\left[(\Lambda_+ + \Lambda_-) - I\right]\\
      &= a^2\left[I - I\right]\\
      &= a^2\left[0\right]\\
      &= 0 \qed
    \end{split}
  \end{equation}

  Let $j = \ket +$.

  \begin{equation}
    \begin{split}
      \prod_{i \ne j}\frac{(A - a_i)}{a_j - a_i} \ket +
      &= \frac{(A - (-a))}{a - (-a)}\ket +\\
      &= \frac{(A + a)}{a + a}\ket +\\
      &= \frac{(a + a)}{a + a}\ket +\\
      &= \ket + \qed\\
    \end{split}
  \end{equation}

  Let $j = \ket -$.

  \begin{equation}
    \begin{split}
      \prod_{i \ne j}\frac{(A - a_i)}{a_j - a_i} \ket -
      &= \frac{(A - a)}{(-a) - a}\ket -\\
      &= \frac{((-a) - a)}{(-a) -a}\ket -\\
      &= \ket - \qed\\
    \end{split}
  \end{equation}

\end{answerpart}

\end{answer}
